\chapter{Panel and time-series comparisons for a banking rollout}
\label{exec17:s6}

A regional rollout of a deposit-servicing program creates several plausible
comparisons. An early region can be compared with a later region, with a
weighted combination of unaffected regions, or with its own forecast under the
old operating rule. Each comparison uses a different source of counterfactual
information. They also answer different questions when treatment effects
vary over the life of a relationship, the donor regions face different funding
shocks, or the bank changes the rule that responds to those shocks.

The earlier panel and time-series chapters establish the identification and
estimation foundations. The task here is to connect the resulting estimands
to executable comparisons without discarding heterogeneity, extrapolation risk
or the calendar on which uncertainty arises. A regression's lead coefficient
may contain post-treatment effects. Exact donor balance may increase mean
squared error. Positive persistence may reduce the variance of an adjacent
difference while leaving a much larger uncertainty for a comparison spanning
several periods. An impulse response under the existing feedback rule need
not describe the response after that rule changes.

We develop these distinctions using exact stipulated laws. The amounts are
illustrative dollars per declared unit of exposure; none is an estimate from
a bank. The constructions supply reproducible mathematical references for a
future public-data or generated-data prototype. Literature claims retain their
specific original versions, while the finite examples and the extensions of
their algebra are project derivations.\citep{exec17_s6_math}

\section{An event-study lead can be made from treatment effects}
\label{exec17:s6-event}

Let $G$ denote the first treatment period, with $G=\infty$ for a never-treated
group. Under an absorbing rollout, the observed outcome in cohort $g$ at time
$t$ has mean
\begin{equation}
 y_{gt}=\mu_{gt}^{0}+1\{t\geq g\}\tau_{gt},
 \qquad \tau_{gt}=E\{Y_t(g)-Y_t(\infty)\mid G=g\}.
 \label{exec17:s6-cohort-law}
\end{equation}
The superscript zero denotes the untreated mean. No anticipation says that
the treated and untreated potential outcomes agree before $g$. It places no
restriction requiring the early and late cohorts to have equal effects after
their own adoption. Sun and Abraham show that pooling relative-time
indicators with unit and time effects can mix treatment effects from other
relative times into a reported coefficient, including a lead. The inspected
original is their September 2020 arXiv version, associated with the 2021
publication; its technical text is not represented as the final journal
version.\citep{sun2021}

Consider an equally weighted, balanced population table with three groups,
$g=2,3,\infty$, and five periods, $t=0,1,2,3,4$. All untreated means are zero.
All treatment effects are zero except $\tau_{2,2}=20$. Thus the early region
has a one-period implementation gain of USD 20, and neither region has any
anticipatory effect. These premises deliberately eliminate a true pretrend.
They isolate a projection problem that would otherwise be confused with a
failure of the rollout's causal assumptions.

For each of the 15 cells, put an intercept, two group indicators and four
time indicators in a nuisance matrix $Z$. Its seven columns are linearly
independent. Put the five event-time indicators
\begin{equation}
 D_{(g,t),k}=1\{g\ne\infty,\ t-g=k\},
 \qquad k\in(-3,-2,0,1,2)
 \label{exec17:s6-design}
\end{equation}
in a matrix $D$. Every event indicator is zero for the never-treated group.
The omitted relative period is $-1$. The extreme lead $-3$ and lag $2$
are included explicitly, rather than silently pooled or dropped. The complete
$15$ by $12$ design $[Z,D]$ has exact rank 12. Consequently the following
coefficients are uniquely defined; no generalized inverse or accidental
collinearity supplies the result.

Residualize against $Z$ by $M_Z=I-Z(Z'Z)^{-1}Z'$. The event coefficients
in the population least-squares projection are
\begin{equation}
 b=(D'M_ZD)^{-1}D'M_Zy.
 \label{exec17:s6-fwl}
\end{equation}
For this example $y$ consists entirely of the five possible post-treatment
cell effects, ordered as
$(\tau_{2,2},\tau_{2,3},\tau_{2,4},\tau_{3,3},\tau_{3,4})'$.
The corresponding columns of the projection operator are
\begin{equation}
 W=
 \begin{pmatrix}
 13/38&4/19&0&-13/38&-4/19\\
 25/171&11/171&0&-25/171&-11/171\\
 11/18&-1/9&0&7/18&1/9\\
 23/114&29/57&0&-23/114&28/57\\
 -8/171&17/171&1&8/171&-17/171
 \end{pmatrix}.
 \label{exec17:s6-projection}
\end{equation}
Rows follow the five event times in \eqref{exec17:s6-design}. These are
projection weights for this declared table, not universal software weights.
For example, the first row subtracts the late cohort's impact and first-lag
effects from positive weights on the early cohort's corresponding effects.
It is not an average of pre-treatment outcomes or pre-treatment causal
effects. The first two row sums are zero, but their individual entries are
nonzero. A zero row sum cannot protect a lead against heterogeneous effects.

Multiplication by $(20,0,0,0,0)'$ yields
\begin{equation}
 (b_{-3},b_{-2},b_0,b_1,b_2)'
 =\left(\frac{130}{19},\frac{500}{171},
          \frac{110}{9},\frac{230}{57},-\frac{160}{171}\right)'.
 \label{exec17:s6-false-leads}
\end{equation}
Both leads are positive although every pre-treatment causal effect is zero.
Even the lag-two coefficient is negative although the only nonzero effect
in the entire table is positive. The example therefore concerns the meaning
of the coefficients before any standard error is computed.

There are two useful checks on this calculation. Direct solution of the full
normal equations gives a residual orthogonal to all 12 columns. Alternatively,
double demean every event indicator and the outcome: subtract its group
mean and time mean and add its overall mean. Regressing the demeaned outcome
on the five demeaned indicators gives exactly
\eqref{exec17:s6-false-leads}. The two calculations use different
representations of the nuisance adjustment. Both retain rational arithmetic,
so a small numerical tolerance cannot hide the false leads.\citep{exec17_s6_math}

The weights explain when cancellation would occur. If the two cohorts had
the same impact effect and the same first-lag effect, the positive and negative
entries in the first two rows would cancel pairwise. The witness violates
that homogeneity while satisfying no anticipation and an identical untreated
law. It does not show that an observed bank pretrend is always spurious.
An actual lead can reflect genuine anticipation, a different untreated trend,
changing composition, or regression contamination. Distinguishing these
possibilities requires a target-preserving estimator and substantive knowledge
of assignment and information.

For a regional program, an apparent gain two periods before adoption should
therefore trigger examination of the comparison weights and rollout clock.
It should not automatically certify early customer response or disqualify the
program. Adding more customers to the same misinterpreted projection can
make its standard error smaller while preserving the wrong population
quantity. In particular, a precise lead test cannot turn
\eqref{exec17:s6-fwl} into an estimate of a different causal target.

The operational alternative is to retain cohort-by-calendar-time contrasts,
their admissible untreated comparisons and their joint uncertainty, then
aggregate to the declared exposure target. The earlier cohort identification
argument remains necessary: choosing a heterogeneous-effect estimator does
not establish parallel trends, support or absence of spillovers.
The inspected \texttt{eventstudyinteract} manual describes an interacted
regression followed by cohort-share weighting, and warns about its control
cohort and balanced-panel validity conditions. Its returned interacted
coefficients and its weighted estimates are distinct objects.
Those restrictions should govern a proposed wrapper, including the choice
of saved covariance; the presence of a familiar event-study plot cannot
serve as semantic admission.\citep{v12_sw_eventstudyinteract}

\section{The horizon determines the trend-sensitivity budget}
\label{exec17:s6-curvature}

Correcting the event-study projection does not settle the counterfactual
trend. We now take as input an admissible treated-minus-comparison mean
series for a fixed population. This can be one cohort or a fixed aggregation
with a common, explicitly defined gap process. It cannot be a sequence whose
cohort composition changes and is then interpreted as the trend of one
unchanging population without further justification.

Index the last untreated period by $-1$, and let $d_t$ be the untreated
treated-minus-comparison gap. Define the last observed untreated slope
$v=d_{-1}-d_{-2}$. For future periods $j=0,\ldots,J-1$, impose the
specific curvature restriction
\begin{equation}
 u_j=d_j-2d_{j-1}+d_{j-2},\qquad |u_j|\leq M.
 \label{exec17:s6-curvature-law}
\end{equation}
The constant $M$ is measured in outcome units per squared calendar period.
It is a maintained sensitivity choice, not the largest observed residual
relabelled as truth. Changing months to quarters changes the restriction.
Rambachan and Roth motivate smoothness and relative-magnitude restrictions
as distinct choices and develop inference that accounts for noisy
pre-treatment information. The construction below specializes a curvature
restriction to one explicit horizon target; it is not an implementation of
their optimized inference procedures.\citep{rambachan2023}

The bias in a baseline-$-1$ contrast at horizon $j$ is
$b_j=d_j-d_{-1}$. Taking first differences in
\eqref{exec17:s6-curvature-law} and summing twice gives
\begin{equation}
 b_j=(j+1)v+\sum_{k=0}^{j}(j-k+1)u_k.
 \label{exec17:s6-triangular}
\end{equation}
Indeed, the next slope is $v+u_0$, the following slope is
$v+u_0+u_1$, and so on. Summing these slopes counts the curvature
$u_k$ exactly $j-k+1$ times. A distant effect is exposed to an early
change of slope repeatedly. This triangular propagation is the reason a
long rollout evaluation can be much more sensitive than its first-period
estimate, even with the same numerical $M$.

Let $\beta_j$ be the observed post-period contrast and $\tau_j$ the
corresponding causal effect, so $\beta_j=\tau_j+b_j$. A fixed target
$\theta=\sum_{j=0}^{J-1}a_j\tau_j$ may be an average effect, a discounted
total or a signed change between horizons. Define
\begin{equation}
 L=\sum_{j=0}^{J-1}a_j(j+1),\qquad
 c_k=\sum_{j=k}^{J-1}a_j(j-k+1).
 \label{exec17:s6-target-kernel}
\end{equation}
Then $\theta=a'\beta-Lv-\sum_k c_ku_k$. Under only the box
restriction in \eqref{exec17:s6-curvature-law}, its sharp identified set is
\begin{equation}
 \Theta(M)=
 \left[a'\beta-Lv-M\sum_k|c_k|,
       \ a'\beta-Lv+M\sum_k|c_k|\right].
 \label{exec17:s6-sharp-curvature}
\end{equation}
To attain the upper endpoint choose
$u_k=-M\operatorname{sign}(c_k)$, with any permitted value when
$c_k=0$. Reverse the signs for the lower endpoint. Starting from the
two fixed pre-period gaps and recursively generating $d_j$ produces a
permitted path for each endpoint. Every intermediate target value is
attainable by convex combinations of the curvature vectors. Sharpness here
uses unrestricted real outcomes and no additional shape constraints; bounded
or otherwise restricted outcomes can narrow the set.

For nonnegative $a_j$, the radius simplifies to
$M\sum_j a_j(j+1)(j+2)/2$. Signed weights require the absolute values
\emph{after} summing the horizon contributions to each curvature. Bounding
each horizon separately and adding absolute target weights is safe but can
be unnecessarily wide. For example, $a=(1,-2,0)'$ has
$c=(-3,-2,0)'$ and radius $5M$, whereas summing the separate absolute
horizon bounds gives $7M$. For $a=(1,-1,1)'$, the kernel is
$(2,1,1)'$ and the radius is $4M$. A package translation must therefore
preserve the target vector, rather than assemble a monetary result from
unrelated marginal sensitivity intervals.

Take pre-period gaps $(-30,-20,-10)$ at times $(-3,-2,-1)$ and
post-period baseline contrasts $(30,35,40)$. These are mean outcome
amounts in dollars per standardized exposure-period. With equal weights
$a=(1,1,1)'/3$, $L=2$, $c=(2,1,1/3)'$, and
\begin{equation}
 a'\beta=35,\qquad Lv=20,\qquad
 \Theta(M)=[15-(10/3)M,\ 15+(10/3)M].
 \label{exec17:s6-money-sensitivity}
\end{equation}
At $M=3$ the effect set is $[5,25]$. The endpoint paths can be
constructed by setting all three future curvatures to $3$ or $-3$.
Exhaustive evaluation of the eight vertices agrees with the triangular
formula. This is a population identification calculation, so its endpoints
do not carry a 95\% coverage label.\citep{exec17_s6_math}

Suppose the recurring incremental operating cost is USD 10 in the same
exposure-period units. Subtracting this fixed cost gives the net-effect set
$[5-(10/3)M,5+(10/3)M]$. A strictly positive lower net value requires
$M<1.5$; at $M=1.5$ it is exactly zero. At $M=3$ the net range is
$[-5,15]$, which leaves the sign unresolved even with the population means
known. This threshold explains what substantive trend restriction an
investment recommendation would require. It does not establish that the
restriction is credible or that the illustrative cost is an actual price.

An average exposure-period benefit must also be distinguished from a
three-period total. With the same three equal exposure counts, the total
has target weights $(1,1,1)'$, center 45 and curvature radius $10M$.
Three periods of the same recurring cost total 30, leaving net center 15.
The positivity threshold remains $M<1.5$, because both benefit and cost
were scaled consistently. Subtracting a single cost of 10 from that total
would answer a different economic question. Unequal exposure counts,
discounting or an upfront fixed cost alter the target weights or cost
functional and require recalculation of the kernel.

\subsection{Retaining the uncertainty in the extrapolated slope}

The preceding center contains an estimated pre-period slope in a real study.
Normalize the gap at $-1$ to zero and let
$\beta_{-2}=d_{-2}-d_{-1}=-v$. Assemble the estimate vector
$\widehat\eta=(\widehat\beta_{-2},\widehat\beta_0,\ldots,
\widehat\beta_{J-1})'$ with its full covariance matrix $\Sigma$.
For fixed target weights, define $h=(L,a')'$. The trend-adjusted center
and its sampling variance are
\begin{equation}
 \widehat m=h'\widehat\eta,\qquad s_m^2=h'\Sigma h.
 \label{exec17:s6-noisy-slope}
\end{equation}
The plus sign on the pre-period coefficient is necessary because the
subtracted slope is $v=-\beta_{-2}$. Even if the post-period contrasts
were estimated independently of one another, the common estimated baseline
would generally induce covariance with the pre-period contrast. Removing the
normalized zero coefficient does not remove that sampling dependence.

For a finite normal design with
$\widehat\eta-\eta\sim N(0,\Sigma)$ and known $\Sigma$, a conservative
$1-\alpha$ interval for the fixed target under the fixed restriction $M$ is
\begin{equation}
 \left[\widehat m-z_{1-\alpha/2}s_m-M\|c\|_1,
       \widehat m+z_{1-\alpha/2}s_m+M\|c\|_1\right].
 \label{exec17:s6-normal-interval}
\end{equation}
On the event $|\widehat m-m|\leq z_{1-\alpha/2}s_m$, every target
permitted by the curvature restriction lies in this interval. That event
has probability $1-\alpha$ under the stipulated Gaussian design. This
proves the stated conservative coverage without identifying the unknown
curvature from a pretest. With an estimated covariance and an asymptotic
normal approximation, a corresponding conclusion needs the approximation
and covariance consistency; finite exact coverage does not follow merely
from inserting a standard error.

For the three-period average, stipulate the diagonal covariance
$\Sigma=\operatorname{diag}(1,4,4,4)$ in squared dollars. This is a
separate hypothetical sampling law, not a covariance inferred from the
15-cell event-study table. Then
\begin{equation}
 h=(2,1/3,1/3,1/3)',\qquad s_m^2=16/3.
 \label{exec17:s6-normal-design}
\end{equation}
Ignoring the noisy slope would retain only variance $4/3$, one quarter
of the correct variance under this particular law. At a 95\% normal
critical value of about 1.96, the sampling radius is about USD 4.53.
For an observed center of 15 and recurring cost 10, a positive lower
net endpoint would require
$5-1.96\sqrt{16/3}-(10/3)M>0$, approximately $M<0.142$.
Population robustness up to $M=1.5$ thus need not imply useful statistical
robustness with this sample precision.

The simple interval is deliberately conservative and uses only the last
pre-period slope. Additional pre-period information, joint restrictions and
better optimized inference can improve efficiency. The comparison with
HonestDiD should retain this distinction. Its documented smoothness and
relative-magnitude procedures use different restriction families; the original
paper discusses why fixed-length, conditional and hybrid approaches have
different domains. A result from one cannot be labelled as another because
their output endpoints happen to resemble each other.\citep{rambachan2023,sw_honestdid}

There is a further selection issue. The proof fixes $M$, the target $a$
and the method before examining favorable post-period estimates. Searching
over horizons or restrictions and reporting the first positive endpoint is
a different random procedure. It requires a joint selection or simultaneous
inference argument. A sensitivity curve can be reported as a collection of
conditional conclusions without claiming that the most favorable point on
the curve has its nominal coverage after selection.

\section{Donor balance is a bias--noise decision}
\label{exec17:s6-donor}

An unaffected donor pool can supply a different counterfactual. Perfect
pre-period fit is attractive, but the future value of extrapolating beyond
that pool depends on noise and contamination as well as balance. To isolate
the tradeoff, take a single post-period untreated outcome law
\begin{equation}
 Y_T(0)=2f+\varepsilon_T,\qquad
 Y_0(0)=\varepsilon_0,\qquad
 Y_1(0)=f+\varepsilon_1.
 \label{exec17:s6-donor-law}
\end{equation}
Here the treated region has known loading 2, the two donors have known
loadings 0 and 1, and the common factor satisfies $|f|\leq B$.
The centered errors are independent with donor variances $\sigma^2$
and treated variance $\sigma_T^2$. The loadings and bound are conditional
premises of the example, not quantities learned exactly from a short panel.

Use donor weights $(1-a,a)$, which sum to one. Their sum removes a common
additive level, if present, but their weighted loading is $a$. The untreated
counterfactual residual is
\begin{equation}
 R(a)=Y_T(0)-(1-a)Y_0(0)-aY_1(0)
      =(2-a)f+\varepsilon_T-(1-a)\varepsilon_0-a\varepsilon_1.
 \label{exec17:s6-donor-residual}
\end{equation}
The same residual is the error in a treatment-effect estimate obtained by
subtracting this donor counterfactual from the treated observation. Conditional
on $f$, its expectation is $(2-a)f$, and its variance is
$\sigma_T^2+\sigma^2\{(1-a)^2+a^2\}$. Taking the worst mean squared
error over the declared factor interval gives
\begin{equation}
 {\cal R}(a)=\sigma_T^2+\sigma^2\{(1-a)^2+a^2\}
                         +B^2(2-a)^2.
 \label{exec17:s6-donor-risk}
\end{equation}
The derivative is
$2\sigma^2(2a-1)+2B^2(a-2)$, so the unique unrestricted minimizer,
when $2\sigma^2+B^2>0$, is
\begin{equation}
 a^*=\frac{\sigma^2+2B^2}{2\sigma^2+B^2}.
 \label{exec17:s6-donor-optimum}
\end{equation}
It lies between $1/2$ and 2. As $B$ approaches zero with positive
noise, the two donors are averaged to reduce variance. As the factor bound
dominates the noise, the weight approaches 2 and exactly matches the
treated loading, at the cost of a negative first-donor weight.

With $B=2$ and $\sigma^2=\sigma_T^2=1$, the complete risk is
\begin{equation}
 {\cal R}(a)=6a^2-18a+18
            =\frac92+6(a-3/2)^2.
 \label{exec17:s6-donor-square}
\end{equation}
Nonnegative weights restrict $a$ to $[0,1]$; the best value on that
interval is $a=1$ with risk 6. Exact loading balance uses $a=2$
and also has risk 6. The intermediate weights $(-1/2,3/2)$ have
risk $9/2$. Their squared bias at the worst factor endpoint is 1,
while their noise variance is $7/2$. The exact fit eliminates this bias
but raises the noise variance to 6. In this law, exact fit spends too much
variance to remove the remaining imbalance.

This calculation is a conditional minimax comparison over a known factor
bound and fixed weights. It is not a general risk guarantee for ridge
augmentation, and $a^*$ is not claimed to be the output of a particular
package. Ben-Michael, Feller and Rothstein develop augmented synthetic control,
including a ridge representation that adjusts balance and can produce
negative weights. Their inspected July 2020 manuscript explains how a penalty
controls departures from the original synthetic-control weights. That
mechanism motivates comparison with the tradeoff above; it does not identify
the bank's $B$ or noise law.\citep{benmichael2021}

Diagnostics should therefore report more than a pre-fit score. The
estimated loading or outcome imbalance, the norm of the weights, and the
validation distribution under which a penalty is selected all matter.
Selecting a penalty on post-treatment outcomes can absorb the effect
being estimated. Selecting it on held-out pre-periods avoids that specific
leakage but still requires an argument that the pre-period prediction task
is informative about the relevant untreated post-period task. A factor that
becomes economically important only after rollout can defeat that transfer.

\subsection{Signed weights also amplify donor contamination}

Suppose the donor observations instead contain unknown mean shifts $c_0,c_1$
with $|c_i|\leq C$. Such shifts may represent spillovers from the program
or an unrecorded intervention in the comparison region. Keeping
\eqref{exec17:s6-donor-law} otherwise fixed, the worst absolute mean error is
\begin{equation}
 B|2-a|+C\{|1-a|+|a|\}.
 \label{exec17:s6-contamination-radius}
\end{equation}
The bound is sharp over the rectangular uncertainty set: choose the sign of
$f$ and each donor shift to align the terms in the residual. If scientific
knowledge links the shifts or restricts their signs, that uncertainty set
should be changed instead of asserting this rectangle is inevitable.

The weight sums are all one, but their absolute sums are 1, 2 and 3
at $a=1,3/2,2$. At a contamination bound $C=3$, the donor component
alone can therefore contribute USD 3, 6 or 9 of bias. Exact loading fit
does not protect against this distinct source of error. In a regional
deposit rollout, transfers of customers or marketing exposure across region
boundaries can be more consequential for a signed counterfactual than for
a nonnegative one, even though both use the same named donors.

If the desired criterion is worst mean squared error with both sources of
bounded bias, the correct objective is
\begin{equation}
 {\cal R}_C(a)=\sigma_T^2+\sigma^2\{(1-a)^2+a^2\}
       +\left[B|2-a|+C\{|1-a|+|a|\}\right]^2.
 \label{exec17:s6-contaminated-risk}
\end{equation}
The two bias budgets add before squaring because an adverse law can align
them. Adding their squared radii would discard the possible cross term.
This distinction is consequential even when every error variance is known.

For $B=2,C=1$ and both variances equal to one, the bias radius is
$5-4a$ for $a\leq0$, $5-2a$ for $0\leq a\leq1$, 3 for
$1\leq a\leq2$, and $4a-5$ for $a\geq2$. The corresponding risks are
\begin{equation}
 {\cal R}_1(a)=
 \begin{cases}
 18a^2-42a+27,&a\leq0,\\
 6a^2-22a+27,&0\leq a\leq1,\\
 2a^2-2a+11,&1\leq a\leq2,\\
 18a^2-42a+27,&a\geq2.
 \end{cases}
 \label{exec17:s6-contaminated-pieces}
\end{equation}
The risk decreases up to the kink at $a=1$ and increases after it.
More explicitly, the stationary points of the four pieces lie outside
their respective intervals, and checking the three finite kinks gives
risks $27,11,15$. The global minimizer is now the nonnegative choice
$a=1$. The previously optimal $a=3/2$ has risk $25/2$, while
exact fit has risk 15.

This reversal is not a demonstration that signed weights should always be
rejected. It demonstrates that a conclusion about regularization depends
on the uncertainty class. A comparison claiming robustness to spillovers
must price that class explicitly. The numerical reference also checks the
three candidate risks by enumerating every endpoint of the factor and
contamination rectangle and every combination of independent centered
unit-variance two-point errors. That direct expectation calculation agrees
with \eqref{exec17:s6-contaminated-risk}.\citep{exec17_s6_math}

\section{Synthetic DiD needs both axes of balance}
\label{exec17:s6-sdid}

The preceding donor example isolates the unit axis. A block rollout also
permits weighting the pre-treatment calendar. To see what this adds, suppose
the untreated mean satisfies
\begin{equation}
 E\{Y_{it}(0)\}=\alpha_i+\gamma_t+\ell_i'f_t,
 \label{exec17:s6-factor-law}
\end{equation}
where $\ell_i$ and $f_t$ are vectors of the same dimension. This is a
conditional factor representation, with additive unit and calendar effects
written separately. We do not infer its rank or predictive validity from a
small pre-fit residual. Let $u_i$ be normalized nonnegative weights on
treated units, $w_i$ normalized donor weights, $v_t$ normalized
nonnegative weights on post periods, and $\lambda_t$ normalized
pre-period weights. Each of the four vectors sums to one over its own
support. Donor and pre-period weights may be signed in the algebra below;
allowing such weights is a comparison exercise, not a description of the
default constraints in an upstream estimator.

Form the fixed-weight double difference
\begin{equation}
 \widehat\tau_{w,\lambda}
 =\sum_{i\in T}u_i\left(\sum_{t\in P}v_tY_{it}
                         -\sum_{t\in B}\lambda_tY_{it}\right)
 -\sum_{i\in C}w_i\left(\sum_{t\in P}v_tY_{it}
                         -\sum_{t\in B}\lambda_tY_{it}\right).
 \label{exec17:s6-double-difference}
\end{equation}
Here $T,C$ are the treated and control sets, while $P,B$ are post and
pre calendars. In a block design, only cells in $T$ by $P$ receive the
treatment. The target causal effect is the treated-post weighted mean
$\tau_{u,v}=\sum_{i\in T,t\in P}u_iv_t\tau_{it}$. This explicit target
allows heterogeneous effects; it should not be silently exchanged for a
simple unweighted ATT when the exposure weights differ.

Normalization cancels the unit effects because the post-minus-pre weights
sum to zero for each unit. It cancels the time effects because the
treated-minus-donor weights sum to zero at each time. The remaining
factor contribution is exactly
\begin{align}
 E(\widehat\tau_{w,\lambda})-\tau_{u,v}
 &=\Delta_\ell'\Delta_f,\label{exec17:s6-imbalance-product}\\
 \Delta_\ell&=\sum_{i\in T}u_i\ell_i-\sum_{i\in C}w_i\ell_i,\nonumber\\
 \Delta_f&=\sum_{t\in P}v_tf_t-\sum_{t\in B}\lambda_tf_t.\nonumber
\end{align}
The proof expands the four rectangles of the factor mean. Every term is
bilinear, so they collect into the product of the two imbalances.
This gives a precise reason to combine unit and time weighting. Eliminating
either imbalance eliminates factor bias under this law. Reducing both can
make their product small without making either exactly zero.

A Euclidean bound is
$|\Delta_\ell'\Delta_f|\leq\|\Delta_\ell\|_2\|\Delta_f\|_2$.
But the separate Euclidean norms depend on how factors are scaled. Replacing
$\ell_i$ by $A\ell_i$ and $f_t$ by $A^{-T}f_t$ for an invertible $A$
leaves every outcome mean and the product unchanged. It need not preserve
the two Euclidean norms. Consequently numerical certificates for separate
loading and factor bounds must declare their normalization; the observed
outcome fit alone does not provide a unique latent scale.

For fixed weights and independent cell errors with common variance
$\sigma^2$, the noise variance of the contrast is
\begin{equation}
 \operatorname{Var}(\widehat\tau_{w,\lambda})
 =\sigma^2(\|u\|_2^2+\|w\|_2^2)
             (\|v\|_2^2+\|\lambda\|_2^2).
 \label{exec17:s6-sdid-noise}
\end{equation}
Each cell's coefficient is the product of its signed unit and time
coefficient. Squaring and summing these products separates into the two
squared norms. This proof requires independent cells and weights fixed
relative to the errors. Serial correlation, shared calendar shocks or
weights fitted to those errors change the variance problem.

For an exact comparison, take one treated unit with loading 2, two donors
with loadings 0 and 1, pre-period factors 0 and 1, and one post-period
factor $3/2$. Begin with donor weights $(0,1)$ and pre-period weights
$(1/4,3/4)$. The loading imbalance is 1 and the factor imbalance is
$3/4$, giving bias $3/4$ and noise variance $(13/4)\sigma^2$.
Changing only the time weights to $(-1/2,3/2)$ matches the post factor
exactly and eliminates the factor bias. It raises the noise variance to
$7\sigma^2$. Changing only the donor weights to $(-1,2)$ also
eliminates factor bias, but gives variance $(39/4)\sigma^2$.
All three calculations can be recovered by multiplying the actual
$3$ by $3$ factor-mean table by its signed unit and time coefficients.

Under these independent-cell premises, their mean squared errors are
\begin{equation}
 {\cal R}_{+}=9/16+(13/4)\sigma^2,\qquad
 {\cal R}_{\mathrm{time}}=7\sigma^2,\qquad
 {\cal R}_{\mathrm{unit}}=(39/4)\sigma^2.
 \label{exec17:s6-sdid-risk}
\end{equation}
The signed time extrapolation improves on the nonnegative comparison
only when $\sigma^2<3/20$. The signed unit comparison has a still
higher variance than the signed time comparison in this particular table.
At $\sigma^2=1$, leaving the small factor bias costs less in mean
squared error than either exact fit. At $\sigma^2=1/10$, signed time
matching has risk $7/10$, below the nonnegative risk $71/80$.
The comparison depends on the post factor, which was stipulated as known
here; an analyst cannot select the favorable matching rule using an
unobserved bank counterfactual.

There is also an important distinction between small mean squared error
and a valid confidence interval. If the remaining bias is not negligible
relative to the standard error, centering a nominal interval at the biased
comparison does not acquire coverage from having lower mean squared error.
A defended factor-bias envelope can instead be carried into a conservative
interval, or a different inferential procedure can be justified for a
specified factor and sampling regime. The risk table is a design
comparison, not an inference theorem.

Arkhangelsky and coauthors develop synthetic DiD with learned unit and
time weights for a block design. The inspected original's formal
normality result assumes independent identically distributed Gaussian
rows with a well-behaved within-row covariance, particular growth
conditions, spectral restrictions and sufficiently predictive, dispersed
oracle weights. It does not require independent calendar cells within a
row, but its independent-row assumption excludes an unrestricted shared
shock across units. The proof also includes a regularization-rate
condition. Neither the fixed-weight identity
\eqref{exec17:s6-imbalance-product} nor a good numerical pre-fit
establishes those premises.\citep{synthdid2021}

The same original distinguishes unit bootstrap, a jackknife that holds
estimated weights fixed, and placebo-based variance estimation. Its
jackknife result uses the theorem's conditions and additional restrictions;
validity is not inferred simply because the estimator has an influence
representation. The single-treated-unit case also requires separate care:
the described jackknife is not defined there, while the placebo
variance argument relies on comparability of untreated noise across
units. This is a substantive reason to retain the inference method in
the output, rather than treating bootstrap, jackknife and placebo as
interchangeable settings.\citep{synthdid2021}

The inspected \texttt{synthdid} documentation requires simultaneous
treatment onset and shows a matrix interface with the number of control
units and pre periods. It also describes the saved interface as beta.
A proposed integration can wrap that declared block route. It should
reject direct admission of staggered adoption, then either choose a
separately justified cohort design or state that the block route does not
answer the request. A mathematical extension discussed in a paper does
not prove that a saved software interface implements it.\citep{sw_synthdid}

\section{Uncertainty belongs to the signed calendar contrast}
\label{exec17:s6-calendar}

A bank can have millions of account rows and few independent regional
funding shocks. At the same time, the direction of serial correlation's
effect depends on the signs in the requested contrast. Consider fixed
unit coefficients $q_i$ and calendar coefficients $r_t$, each summing
to zero. A double difference has the scalar form
$\widehat\theta=\sum_iq_i\sum_tr_tY_{it}$. Suppose its untreated
noise is
\begin{equation}
 e_{it}=\ell_i b_t+\epsilon_{it},\qquad
 \operatorname{Cov}(b_t,b_s)=\Gamma_{ts},
 \label{exec17:s6-shared-law}
\end{equation}
with independent, centered cell errors $\epsilon_{it}$ of common
variance $\sigma^2$, independent of $b$. The loadings $\ell_i$ are
fixed. Let $\kappa=\sum_iq_i\ell_i$. The shared component of the
contrast is $\kappa\sum_tr_tb_t$, so
\begin{equation}
 \operatorname{Var}(\widehat\theta)
   =\kappa^2 r'\Gamma r+\sigma^2\|q\|_2^2\|r\|_2^2.
 \label{exec17:s6-signed-variance}
\end{equation}
This is a finite-calendar covariance calculation, with no large-sample
claim hidden in the notation. If the unit contrast balances the common
loading exactly, $\kappa=0$, this particular common shock cancels.
The normalization $\sum_iq_i=0$ cancels only equal loadings; it does
not imply $\kappa=0$ when regions have different exposure.

Let $b_t$ be stationary AR(1), with marginal variance $\sigma_b^2$
and correlation $\rho^{|t-s|}$, $|\rho|<1$. For the gap-$h$
difference $r=(-1,0,\ldots,0,1)'$,
\begin{equation}
 r'\Gamma r
  =\operatorname{Var}(b_h-b_0)
  =2\sigma_b^2(1-\rho^h).
 \label{exec17:s6-gap-variance}
\end{equation}
The covariance enters with a minus sign. For a short difference and
positive persistence, nearby common shocks partially cancel. For a
longer gap they cancel less. This observation does not support ignoring
serial correlation; it shows why transferring a variance-of-a-level-mean
intuition to a signed treatment contrast can be wrong in either direction.

An independent check starts from
$b_h=\rho^hb_0+\sum_{j=1}^h\rho^{h-j}\eta_j$, with innovations
independent of $b_0$ and variance
$\sigma_b^2(1-\rho^2)$. Then
\begin{align}
 \operatorname{Var}(b_h-b_0)
 &=\sigma_b^2(\rho^h-1)^2
   +\sigma_b^2(1-\rho^2)\sum_{j=0}^{h-1}\rho^{2j}\nonumber\\
 &=2\sigma_b^2(1-\rho^h).
 \label{exec17:s6-innovation-check}
\end{align}
This derivation verifies the finite covariance result without constructing
a Toeplitz covariance matrix. It also exposes why the stationary
initial-state variance matters. Starting the process at a known fixed
zero state would give a different finite-calendar law.

For one treated region and one donor, set $q=(1,-1)'$ and
$\kappa=1$. With $\sigma_b^2=100$, $\rho=4/5$ and $\sigma^2=1$,
the adjacent difference has common variance 40 and idiosyncratic variance
4, for total 44. The gap-three difference has common variance
$488/5=97.6$ and total $508/5=101.6$. Both use the same two
regions and the same number of observed endpoints. Their uncertainty
differs because of the interval they span, not because one has more
customer records.\citep{exec17_s6_math}

To examine customer replication explicitly, suppose each region-period
outcome is the mean of $n$ customers with independent unit-variance
idiosyncratic errors and the same regional loading. The idiosyncratic
variance of that mean is $1/n$, so the two total variances become
$40+4/n$ and $97.6+4/n$. As $n$ increases, the limits are 40
and 97.6. This calculation gives an economic reason to acquire more
independent regions or periods when feasible: an additional customer in
an already observed shock cell cannot remove the common component.
It does not prove that any given additional region is independent,
uncontaminated or comparable.

The covariance also interacts with weighting. In the double difference,
$q$ contains the treated and signed donor weights, and $r$ contains
post and signed pre weights. Choosing a weight to balance a factor can
reduce $\kappa$ while increasing $\|q\|_2$ or $\|r\|_2$.
The independent-cell variance formula
\eqref{exec17:s6-sdid-noise} is therefore a special case of a more
specific noise law, rather than a generic cost attached to every
synthetic comparison. If weights are fitted from the same outcome
errors, their randomness adds a further issue beyond the covariance
of a fixed contrast.

An inferential route must specify which dimension supplies independent
replication and what remains conditional. Resampling customers does not
reproduce a region-level common shock. Resampling whole units respects
within-unit trajectories but does not automatically reproduce shocks
shared across all units. A time-block route needs a defended dependence
regime and enough calendar information; a short fixed calendar cannot
be made asymptotically long by changing the resampling option. The
earlier covariance and dependence chapters provide possible methods.
Here the requirement is to evaluate them for the actual signed target,
including its loading imbalance and its information clock.

\section{A response under feedback is a policy-specific response}
\label{exec17:s6-feedback}

A regional time series supplies yet another comparison: forecast the outcome
under the existing operating rule and perturb an action. The intervention must
be specified together with the rule that will react to its consequences.
Consider the structural system
\begin{equation}
 Y_t=\rho Y_{t-1}+\beta A_t+\epsilon_t,\qquad
 A_t=kY_{t-1}+u_t.
 \label{exec17:s6-feedback-law}
\end{equation}
The state $Y_t$ might be a standardized funding position, and $A_t$ a
signed adjustment in servicing or funding intensity. The parameters and
innovation interpretation are supplied structural premises. A fitted
autoregression does not by itself establish that its residual is an
exogenous, manipulable action shock. We compare differences between paths
with the same initial state and structural innovations, so the additive
$\epsilon_t$ cancels in the response calculation.

Introduce a unit action innovation at time zero, $u_0=1$, and set later
innovations to zero. If the operating feedback remains $k$, substitution
into the state equation gives the response
\begin{equation}
 R_h^{\mathrm{maint}}=\beta(\rho+\beta k)^h,
 \qquad h=0,1,\ldots.
 \label{exec17:s6-maintained-response}
\end{equation}
The first state response is $\beta$. Its effect on the next action is
$k\beta$, which adds $\beta k\beta$ to the next state response in
addition to its inherited persistence $\rho\beta$. Repeating this
recursion proves the power in \eqref{exec17:s6-maintained-response}.
The later action responses are $k\beta(\rho+\beta k)^{h-1}$ for
$h\geq1$. A unit innovation is therefore followed by a sequence of
endogenous adjustments; it is not a one-period action path.

If the bank instead forces $A_0=1$ and every later action difference to
zero, the response is
\begin{equation}
 R_h^{\mathrm{open}}=\beta\rho^h.
 \label{exec17:s6-open-response}
\end{equation}
If it replaces the feedback coefficient by $k'$ and then introduces the
same unit innovation, the corresponding response is
\begin{equation}
 R_h^{\mathrm{replace}}=\beta(\rho+\beta k')^h.
 \label{exec17:s6-replacement-response}
\end{equation}
This last calculation assumes that the structural outcome law and innovation
meaning survive the replacement. A rule change that also changes the outcome
mechanism needs a different structural analysis. Moreover, it is the
incremental response to an innovation \emph{within} the replacement regime;
it is not the entire value of switching regimes from an arbitrary nonzero
initial portfolio.

Take $\rho=1/2$, $\beta=1$, $k=2/5$ and $k'=-1/5$. At horizons
zero through three the three state paths are
\begin{align}
 R^{\mathrm{maint}}&=(1,0.9,0.81,0.729),\nonumber\\
 R^{\mathrm{open}}&=(1,0.5,0.25,0.125),\nonumber\\
 R^{\mathrm{replace}}&=(1,0.3,0.09,0.027).
 \label{exec17:s6-feedback-paths}
\end{align}
All three have the same impact response. Their four-period sums are
$3.439$, $1.875$ and $1.417$. At USD 100 of margin per state-period
and an upfront USD 200 implementation cost, the resulting net amounts,
before charging later action intensity, are USD 143.90, $-12.50$ and
$-58.30$. A positive maintained-rule calculation can consequently coexist
with a negative value for the forced one-period action path. The discrepancy
is a difference between interventions, rather than an estimator's rounding
error.\citep{exec17_s6_math}

For this numerical example the action can be positive or negative without a
capacity bound. A real servicing action may be nonnegative, integer-valued or
capped. Enforcing those restrictions changes the recursion when they bind.
Neither the powers above nor a linear software response should then be
presented as the response of the constrained policy without solving that
policy's state evolution.

\subsection{Closed-loop fit can leave the replacement response unidentified}

Even a perfect fit to the old feedback system need not identify these
interventions. Suppose the observed action is always $A_t=(2/5)Y_{t-1}$,
with no action innovation. Consider two structural worlds:
\begin{equation}
 (\rho,\beta)=(1/2,1)\quad\hbox{or}\quad(1/10,2).
 \label{exec17:s6-two-worlds}
\end{equation}
Both imply $Y_t=(9/10)Y_{t-1}+\epsilon_t$ under the observed rule.
With the same initial law and innovation law they generate the same joint
distribution of all observed states and actions. Yet replacing the feedback
by $k'=-1/5$ gives state coefficients $3/10$ and $-3/10$ respectively.
The observed reduced-form persistence and the known old action rule do not
separate the state persistence from the action effect.

For a known initial state $Y_{-1}=1$ and zero subsequent innovations, the
four replacement-rule states are $(0.3,0.09,0.027,0.0081)$ in the first
world and $(-0.3,0.09,-0.027,0.0081)$ in the second. Their sums are
$0.4251$ and $-0.2289$, or USD 42.51 and $-22.89$ of gross margin at
the same conversion. These are state paths after a controller replacement,
whereas \eqref{exec17:s6-feedback-paths} describes innovation responses
within each regime. Keeping that distinction prevents another target
substitution.

The witness is deliberately a lack-of-excitation case. Exogenous variation
in actions, a defended instrument, a structural restriction on $\rho$ or
$\beta$, or a new experiment can provide additional information. A high
forecast score under the old controller cannot supply it. An estimator
request should therefore carry both the action-support evidence and the
intervention being asked of the time series. When these are absent, a
set-valued answer or refusal is more faithful to the observed information
than a unique replacement-policy response.

\section{Value includes the action path and the unobserved tail}
\label{exec17:s6-value}

The same horizon discipline used for the trend budget applies to feedback
value. For a deterministic zero-baseline example, let $m_h$ be dollars per
state unit, $d\in(0,1]$ a per-period discount factor, $q\geq0$ the price
of squared action intensity, and $K$ an upfront fixed cost. The finite
incremental value is
\begin{equation}
 V_H=\sum_{h=0}^{H}d^h\{m_hR_h-q(\Delta A_h)^2\}-K.
 \label{exec17:s6-finite-value}
\end{equation}
This cost is a stipulated capacity penalty, not a market price. Its
quadratic form is useful for exposing the extra future actions hidden in
a maintained-feedback response. At a nonzero baseline action $A_h^0$,
the change in a quadratic cost is instead
$q\{2A_h^0\Delta A_h+(\Delta A_h)^2\}$. Ignoring the cross term
would generally misprice the incremental policy. Likewise, uncertainty
about actual action paths cannot be priced by squaring only their mean.

For the zero-baseline paths above, the four action responses are
\begin{align}
 \Delta A^{\mathrm{maint}}&=(1,0.4,0.36,0.324),\nonumber\\
 \Delta A^{\mathrm{open}}&=(1,0,0,0),\nonumber\\
 \Delta A^{\mathrm{replace}}&=(1,-0.2,-0.06,-0.018).
 \label{exec17:s6-action-paths}
\end{align}
At $d=1$, $q=10$, $m_h=100$ and $K=200$, their squared action sums
are $1.394576$, $1$ and $1.043924$. Subtracting the corresponding
penalties changes the four-period net values to USD 129.95424,
$-22.50$ and $-68.73924$. The values follow both from summing the
displayed paths and from a separate structural state/action recursion.
They establish an arithmetic comparison under the declared law, not the
economic superiority of a bank policy.\citep{exec17_s6_math}

For a stable response $R_h=\beta a^h$ with maintained coefficient $k$,
constant margin $m$ and $H\geq0$, geometric sums give
\begin{align}
 V_H={}&m\beta\frac{1-(da)^{H+1}}{1-da}-K\nonumber\\
 &-q\left[1+k^2\beta^2d\,
                   \frac{1-(da^2)^H}{1-da^2}\right].
 \label{exec17:s6-geometric-value}
\end{align}
Here the impact action is one, and $\Delta A_h=k\beta a^{h-1}$
thereafter. The expression is interpreted by continuity when a denominator
vanishes, using the corresponding finite sum. Setting $k=0$ and
$a=\rho$ recovers the open-loop action path; setting $k=k'$ and
$a=\rho+\beta k'$ gives the replacement-regime innovation calculation.
Thus a common formula can compare the three interventions without
quietly carrying the old feedback's actions into the new one.

If $|da|<1$ and $da^2<1$, the infinite discounted value is
\begin{equation}
 V_\infty=\frac{m\beta}{1-da}-K
       -q\left[1+\frac{k^2\beta^2d}{1-da^2}\right].
 \label{exec17:s6-infinite-value}
\end{equation}
These are different convergence requirements. Discounting a state response
can make its benefit series converge while the squared action cost series
still diverges. For example, $d=1/2$ and $a=3/2$ give $da=3/4$ but
$da^2=9/8$. With positive $q$ and nonzero $k\beta$, finite discounted
benefit then coexists with an infinite action cost in this unconstrained
law. A reporting engine should not declare a finite net present value
after checking only the response's first-moment discount condition.

At $d=0.95$, the maintained, open and replacement example gives infinite
net values of approximately USD 473.06, $-19.52$ and $-70.56$.
These calculations use the same known parameters forever. Their large
difference from four-period values is a warning about the economic weight
placed on persistence, not evidence that the far future is predictable.
A structural break, capacity constraint or an unstable parameter in a
confidence region invalidates the corresponding plug-in extrapolation.

\subsection{A conditional truncation budget}

An infinite-horizon claim needs a defended tail restriction. Suppose
uniformly over a declared admissible regime
$|m_hR_h|\leq C r^h$ with known $C\geq0$, $r\geq0$ and $dr<1$.
Then the absolute omitted gross benefit after horizon $H$ is bounded by
\begin{equation}
 T_B(H)=C\frac{(dr)^{H+1}}{1-dr}.
 \label{exec17:s6-tail-budget}
\end{equation}
This follows by applying the triangle inequality to each omitted term
and summing the geometric series. It is an approximation budget
conditional on a uniform structural envelope, rather than a sampling
standard error. With $C=100$, $d=0.95$, $r=0.9$ and $H=3$, it is
about USD 368.55. Reporting the four observed periods to cents does
not reduce this stipulated omitted-benefit budget.

For the linear action envelope
$|\Delta A_h|\leq |k\beta|r^{h-1}$ at $h\geq1$, the omitted
quadratic cost after $H\geq0$ is nonnegative and bounded, if $dr^2<1$,
by
\begin{equation}
 T_A(H)=qk^2\beta^2\frac{d^{H+1}r^{2H}}{1-dr^2}.
 \label{exec17:s6-action-tail}
\end{equation}
Consequently, absent further sign information on the margin response,
the omitted net value belongs to
$[-T_B(H)-T_A(H),\ T_B(H)]$. This uses the known nonnegative
sign of the cost instead of adding a symmetric cost allowance to both
ends. If benefits are also known nonnegative, the lower endpoint can
be tightened to $-T_A(H)$. Such sign information is another premise;
it is not supplied by the word ``benefit.''

If $C=0$, the gross tail is zero at every horizon. For $C>0$, an ex ante
gross-tail tolerance $\varepsilon>0$ and $0<dr<1$,
the smallest admissible nonnegative integer horizon satisfies
\begin{equation}
 H+1\ \geq
 \frac{\log\{\varepsilon(1-dr)/C\}}{\log(dr)}.
 \label{exec17:s6-horizon-budget}
\end{equation}
The reversed inequality follows because $\log(dr)<0$. In the numerical
envelope, a USD 10 gross-tail tolerance needs $H=27$: $H=26$
still has a bound above USD 10, whereas $H=27$ is below it.
This is a requirement for the chosen envelope, not an instruction to
forecast 28 periods (indices $0$ through $27$) from an uninformative short history. More credible
long-horizon restrictions, a finite contractual decision horizon, or
an explicit unresolved value set may be preferable.

In statistical use, a simultaneous confidence region for structural
parameters can be propagated through finite-horizon value and a tail
bound valid uniformly on that region. If the region includes a parameter
with an unbounded adverse tail, its corresponding lower value cannot
be replaced by the value at the stable point estimate. A claim about
the region's coverage remains conditional on identification, dependence
and estimation assumptions; the deterministic tail inequality supplies
none of them. This separation makes it possible to report parameter,
trend, donor, numerical and horizon limitations in their correct units
without calling every budget a confidence interval.

\section{Software routes must preserve the comparison}
\label{exec17:s6-software}

The proposal should reuse a mature estimator where its target and
sampling conditions fit. The relevant choice is more specific than
``run a panel package.'' The required state includes the rollout
calendar, the permitted untreated comparisons, the population and
horizon weights, and the observation or innovation law. The selected
primary documentation supports three complementary routes below.
It is a bounded comparison of inspected versions; the wider software
survey, complete dependency rights and runtime qualification remain
separate obligations.\citep{exec17_s6_paths}

\subsection{Cohort estimates followed by trend sensitivity}

The first route uses a cohort-aware event-study estimator and passes
a correctly normalized estimate vector with its joint covariance to
a trend-sensitivity procedure. The inspected HonestDiD example for
\texttt{did} requires a dynamic event study, a universal base period,
consecutive relative times and a reference period of $-1$. It
constructs covariance from the dynamic influence vectors and removes
the normalized zero coefficient. Its example target vector selects
one post period. A discounted banking total needs an explicitly
changed target vector and matching outcome units, rather than adding
the endpoints of several separately reported intervals.
\citep{sw_honestdid,exec17_s6_paths}

The bridge's influence-product covariance is not a universal remedy
for shared regional calendar shocks. Its validity depends on the
sampling units represented by the input influence rows and the
underlying estimator's inferential regime. A proposed wrapper must
either justify the joint covariance for the bank's actual design or
decline the requested inferential claim. Retaining the whole matrix
also permits the noisy-slope contrast in
\eqref{exec17:s6-noisy-slope}; exporting only marginal standard
errors would lose that possibility.

The selected \texttt{eventstudyinteract} route is a community Stata
estimator. Its saved \texttt{e(b\_iw)} and \texttt{e(V\_iw)} describe
the weighted estimates; \texttt{e(b)} and \texttt{e(V)} retain the
underlying interacted regression. A subprocess wrapper should export
the intended pair, its cohort weights and support, rather than rely
on whichever object a generic post-estimation command reads. The
manual says to exclude periods after the last cohort becomes treated
when that cohort supplies the controls. It also warns that the cited
validity result is for balanced panels without covariates, although
the implementation accepts unbalanced data. Admitting the broader
input format does not establish the broader causal or inferential
claim.\citep{v12_sw_eventstudyinteract}

We propose wrapping these existing routes when their regimes and
rights are satisfied. Reimplementing their inference in another
language has no demonstrated benefit at this stage and would create
new numerical and semantic obligations. Stata runtime rights and
the community package's unresolved redistribution rights must also
remain distinct. Their absence from this bounded inspection is a
reason to defer redistribution, not evidence that the estimator is
scientifically inferior.

\subsection{Augmented donors and block synthetic DiD}

The inspected \texttt{augsynth} dispatcher has separate single-outcome,
multiple-outcome and staggered paths. It exposes outcome-model,
balance and regularization choices, with some options specific to
the dispatched path. The saved declaration is version 0.2.0 at its
recorded source commit. We propose evaluating the original R path
and preserving the dispatcher result, donor identities, penalty,
weight signs and chosen uncertainty method. The entry-point manual
does not establish every downstream implementation or inferential
detail; those need inspection before admission. In particular, the
oracle minimizer in \eqref{exec17:s6-donor-optimum} should not be
used as a claim about the output of its ridge option.
\citep{v12_sw_augsynth,exec17_s6_paths}

The proposed \texttt{synthdid} route is narrower: the inspected
matrix interface uses an untreated-unit count and a pre-period
count with simultaneous treatment onset. The transformation from
long banking observations to that matrix must preserve the meaning
of both axes. Dropping early-treated cells to create a rectangle
changes the observed comparison and may change the population
target. For a staggered request, a collection of separately justified
cohort comparisons is an alternative design, rather than a hidden
preprocessing trick that makes the block estimator universal.
Its beta-interface status also argues for a version-bound wrapper
before considering a port.\citep{sw_synthdid}

These routes solve related but different balance problems. Augmentation
can spend extrapolation to improve unit balance; synthetic DiD also
changes the pre-period comparison. The product identity
\eqref{exec17:s6-imbalance-product}, norm penalty
\eqref{exec17:s6-sdid-noise} and contamination radius
\eqref{exec17:s6-contamination-radius} explain concrete reasons their
outputs can differ. They do not provide a ranking until the uncertainty
class, sampling law, common target and actual estimator behavior are
specified. We therefore propose retaining both original engines for
matched evaluation. We reject automatic substitution between their
outputs and defer a native port until a measured requirement justifies
its maintenance and verification cost.

\subsection{A prototype can expose semantic errors before a bank trial}
\label{exec17:s6-prototype}

A useful first prototype compares meanings before comparing leaderboard
scores. The exact event-study table supplies a known false-lead
reference, including the full design rank and projection matrix.
The curvature laws supply sharp endpoints and a noisy-slope variance
check. The donor and double-weight tables supply known bias, variance
and contamination tradeoffs. The signed-calendar and feedback laws
then expose two failures that a purely tabular fit score can miss:
the wrong source of replication and the wrong intervention after
a rule change. These are proposed admission tasks, supported by
local authored algebra; no upstream estimator has been run here.
\citep{exec17_s6_math,exec17_s6_design}

Permitted public panels can make the exercise more realistic by supplying
calendar spacing, exposure sizes, missingness and covariate patterns.
Their license, revision date, units and local identity would be retained.
They do not by themselves reveal the untreated outcome of a treated
region. One defensible design overlays explicitly generated potential
outcomes on a permitted public covariate structure, disclosing that
the causal truth is constructed. Another compares descriptive behavior
on a real policy episode while preserving unresolved causal truth.
Those designs answer different questions and should be reported
separately. Existing public-data survey entries can be screened for
these purposes without treating a dataset's availability as scientific
validation of the rollout design.

After semantic admission, a prospective statistical comparison would
freeze each engine and its dependency environment, the common
estimand, a fresh generating law and seed, sample sizes, repetitions,
inference methods, stopping rule and failure accounting. Scenarios
would vary treatment-effect heterogeneity, factor extrapolation,
donor contamination, calendar dependence, action excitation and
post-change stability. Methods should be compared only where their
targets coincide and their claimed regimes are declared; a justified
refusal is an observable outcome. No favorable preliminary result
should be used to revise the final scenario denominator.

The resulting report would retain target error, interval coverage
under the constructed law, availability, width, monetary decision
loss and observed end-to-end cost. It would also retain unsupported
requests, convergence failures, nonfinite outputs and unresolved
rights. Statistical qualification of one frozen regime would remain
different from acceptance of arbitrary panels or a bank's actual
assignment and shock law. Only after such a prototype demonstrates
useful behavior is a costly authentic-bank trial justified by evidence
from this program.\citep{exec17_s6_design}

For the proposed library, CLI and MCP interface, a single result should
therefore carry the same target vector, reference period, supported
cells, signed weights, uncertainty method, sensitivity restrictions
and horizon budget. A scalar ``effect'' field would erase the
distinctions developed here. The transport should also preserve a
set-valued result or refusal when assumptions or support do not
admit a scalar claim. This is a proposed semantic contract for later
integration; the documented estimator interfaces and these exact
references do not establish interoperability with another repository.

\section{Conclusion}
\label{exec17:s6-end}

A panel or time series supplies a comparison only after the target,
counterfactual information and sampling law are specified. The exact
rollout regression produces positive leads without anticipation.
A target-specific trend budget becomes more demanding with the
horizon and with uncertainty in the pre-period slope. Exact donor
balance can lose to a regularized comparison, while an explicit
contamination class can reverse that ranking. Synthetic DiD's two
axes of balance interact with noise and shared shocks. A maintained
feedback response can value a different intervention from the one
the bank proposes to implement.

These results motivate concrete reuse and prototype routes. They
also identify the information still needed: credible counterfactual
restrictions, appropriate independent replication, action excitation,
stable intervention meaning, feasible costs and a defended horizon.
An integrated proposal should preserve those requirements through
the estimator, uncertainty calculation and monetary decision. A
successful software call cannot replace them.
